% Owner: TV2 (P5.2a) -- Reviewer: TV3
\subsection{Choosing the Final Move by $N$ or by $W/N$}
\label{subsec:discussion_n_vs_q}

A move with a high mean $W/N$ but few visits is a move whose estimate is still uncertain: with $N_i$ simulations, the confidence radius of~\eqref{eq:radius} shrinks only as $1/\sqrt{N_i}$, so a child that won $2$ of its $2$ simulations has a perfect mean and a very wide interval. UCT already accounts for this during the search. A child keeps receiving visits only while its upper bound remains competitive, and by~\eqref{eq:plays} a clearly worse move receives only $O(\ln n)$ visits. The visit count therefore summarizes both the estimated value and the confidence in it, which is why the most visited child is the standard choice~\cite{browne2012survey}.

Our results support this view. On the trap positions of Table~\ref{tab:t4}, the three rules agree whenever the budget is large enough for the correct move to dominate the visits, as on E1, and they disagree when the statistics at the root are still noisy, as on the diagonal fork and the $9\times 9$ open three. In actual games (Table~\ref{tab:t5}), only $2$ of $27$ comparisons are significant, in opposite directions. The max-robust rule, which falls back on a lower confidence bound when $N$ and $W/N$ disagree, was never significantly worse than the others. We therefore use the most visited child by default. It costs nothing, and it agrees with the theory. Max-robust remains a reasonable alternative when the budget is small. The more important lesson of Table~\ref{tab:t4} is that no selection rule can recover a move that the search did not explore: on the $9\times 9$ open-three position all rules fail in about $87\%$ of the runs, and the remedy lies in the simulations and in prior knowledge (Tables~\ref{tab:t6} and~\ref{tab:b2b3}), not in the final choice.
